\section{Introdução}

\subsection{Adendos}

\begin{frame}{Antes de começarmos....}
    \begin{itemize}[<+-| alert@+>] 
      % Bom pq são uteis, varias definições e relações com as artes das trevas
      \item Aparecem em todo canto
        \item Não são uma classe de grafo (meio que sim mas
          não)
    \end{itemize}
\end{frame}


\subsection{Motivação}

\begin{frame}[c]{Imagine...}

  Imagine a situação cotiadiana em que você precisa de um grafo com N vértices e
  com poucas arestas, e que ao passear sobre ele dá pra chegar em qualquer lugar rapidinho.
  \vspace{40pt}
  \onslide<2->{%
    \begin{center}
      \textcolor{red}{\Huge O que você faria?}
    \end{center}
  }
\end{frame}

\begin{frame}[fragile]{Hmmm...}

  Talvez um grafo completo...
  \begin{center}
    \begin{tikzpicture}[scale=1.5]
      \node (K5) {
        \begin{tikzpicture}
          \graph [nodes={circle,draw}, clockwise] {
            subgraph K_n [n=5];
          };
        \end{tikzpicture}
      };

    \node[below=10pt] at (K5.south) {$n=5\text{ }m=10$};
    \visible<2->{%
      \node[right] (K6) at (K5.east) {
        \begin{tikzpicture}
          \graph [nodes={circle,draw}, clockwise] {
            subgraph K_n [n=6];
          };
        \end{tikzpicture}
      };
    \node[below=10pt] at (K6.south) {$n=6\text{ }m=15$};
    };

    \visible<3->{%
      \node[right] (K7) at (K6.east) {
        \begin{tikzpicture}
          \graph [nodes={circle,draw}, clockwise] {
            subgraph K_n [n=7];
          };
        \end{tikzpicture}
      };
    \node[below=10pt] at (K7.south) {$n=7\text{ }m=21$};
    };

    \end{tikzpicture}   

    \vspace{20pt}
    \onslide<4->{%
      Puts!! Tem arestas $\mathcal{O}(N^2)$.
    }
  \end{center}
\end{frame}

\begin{frame}{Não se preocupe!!!}
  Temos a solução para suas necessidades:
  \vspace{20pt}
  \begin{center}
    \textcolor{red}{\Huge Os Grafos Expansores}
  \end{center}
\end{frame}


\begin{frame}{Mas antes....}
  \vspace{20pt}
  \begin{center}
  Precisamos conversar sobre "chegar rapidinho em qualquer lugar"
  \end{center}
\end{frame}
